\Askhsh{17956}{4}
\begin{ylh}{1}
    \item 10.5. Λόγος εμβαδών όμοιων τριγώνων - πολυγώνων.
\end{ylh}
Δίνεται τρίγωνο $\text{ΑΒΓ}$ και σημείο $\text{Δ}$ στο εσωτερικό του τμήματος $\text{ΒΓ}$. Από το $\text{Δ}$ φέρουμε παράλληλες στις πλευρές $\text{ΑΒ}$ και $\text{ΑΓ}$. Η παράλληλη στην $\text{ΑΒ}$ τέμνει την $\text{ΑΓ}$ στο $\text{Ζ}$ και η παράλληλη στην $\text{ΑΓ}$ τέμνει την $\text{ΑΒ}$ στο $\text{Ε}$. Θεωρούμε $\text{Κ}$ και $\text{Λ}$ τα μέσα των $\text{ΒΔ}$ και $\text{ΔΓ}$ αντίστοιχα. Να αποδείξετε ότι:
\begin{alist}
\item $(\text{ΕΚΔ}) = \frac{(\text{ΒΕΔ})}{2}$ \monades{09}
\item $(\text{ΕΖΔ}) = \frac{(\text{ΑΕΔΖ})}{2}$ \monades{09}
\item Το εμβαδόν του $\text{ΚΕΖΛ}$ είναι ανεξάρτητο της επιλογής του σημείου $\text{Δ}$. \monades{07}
\end{alist}